\section*{Bab \ref{ch:b2:plant-life-cycles} --- Daur Hidup dan Perkembangbiakan Tumbuhan Darat}

\begin{solution}{exo:b2:plant-life-cycles:1}
\omterm{def:b2:plant-life-cycles:cycles}{Gametofit}: keturunan \omterm{def:b2:meiosis-heredity:meiosis}{haploid} \omterm{def:b2:unicellular-diversity:unicellular}{bersel banyak}, yang tumbuh secara mitosis
dari sebuah spora, dan yang membuat gamet secara mitosis. \omterm{def:b2:plant-life-cycles:cycles}{Sporofit}:
keturunan \omterm{def:b2:meiosis-heredity:meiosis}{diploid}, yang tumbuh secara mitosis dari sebuah zigot, dan
yang membuat spora secara \omterm{def:b2:meiosis-heredity:meiosis}{meiosis}. Spora: sel \omterm{def:b2:meiosis-heredity:meiosis}{haploid} yang dibuat
secara \omterm{def:b2:meiosis-heredity:meiosis}{meiosis} dan tumbuh tanpa melebur. Gamet: sel \omterm{def:b2:meiosis-heredity:meiosis}{haploid} yang dibuat
secara mitosis (pada tumbuhan) dan harus melebur dengan gamet lain.
\end{solution}

\begin{solution}{exo:b2:plant-life-cycles:2}
Sel daunnya $2n = 1260$: spora 630, sel \omterm{def:b2:plant-life-cycles:fern}{protalus} 630, sperma 630, sel
telur 630, zigot 1260.
\end{solution}

\begin{solution}{exo:b2:plant-life-cycles:3}
Bantalan hijaunya: \omterm{def:b2:plant-life-cycles:cycles}{gametofit}. Tangkai dan kapsulnya: \omterm{def:b2:plant-life-cycles:cycles}{sporofit}. Ental
pakunya: \omterm{def:b2:plant-life-cycles:cycles}{sporofit}. \omterm{def:b2:plant-life-cycles:fern}{Protalusnya}: \omterm{def:b2:plant-life-cycles:cycles}{gametofit}. Butir \omterm{def:b2:plant-life-cycles:heterospory}{serbuk sarinya}:
\omterm{def:b2:plant-life-cycles:cycles}{gametofit} jantan. Simpanan makanan bijinya (pada pinus): \omterm{def:b2:plant-life-cycles:cycles}{gametofit}
betina. Lembaga bijinya: \omterm{def:b2:plant-life-cycles:cycles}{sporofit} berikutnya.
\end{solution}

\begin{solution}{exo:b2:plant-life-cycles:4}
Haplontik (\emph{Chlamydomonas}): \omterm{def:b2:meiosis-heredity:meiosis}{meiosis} di dalam zigotnya, seketika.
Diplontik (hewan, \emph{Fucus}): \omterm{def:b2:meiosis-heredity:meiosis}{meiosis} membuat gametnya.
Haplo-diplontik (\omterm{def:b2:plant-life-cycles:moss}{lumut}, \omterm{def:b2:plant-life-cycles:fern}{tumbuhan paku}, tumbuhan berbiji, banyak alga):
\omterm{def:b2:meiosis-heredity:meiosis}{meiosis} membuat spora di dalam \omterm{def:b2:plant-life-cycles:fern}{sporangium} \omterm{def:b2:plant-life-cycles:cycles}{sporofitnya}.
\end{solution}

\begin{solution}{exo:b2:plant-life-cycles:5}
$0.03/10^{-4} = \qty{300}{s}$, yaitu lima menit. Dalam \qty{20}{min},
sejauh \qty{12}{cm}. Pembuahannya hanya berhasil antara tumbuhan yang
berjarak beberapa sentimeter, sehingga \omterm{def:b2:plant-life-cycles:moss}{lumut} tumbuh dalam bantalan
rapat dengan kedua kelaminnya (atau kedua organnya) dalam jangkauan.
\end{solution}

\begin{solution}{exo:b2:plant-life-cycles:6}
$v_s = 2\times(1.5\times 10^{-5})^2\times 1099\times 9.81/(9\times
1.8\times 10^{-5}) = \qty{0.030}{m/s}$, yaitu \qty{3}{cm/s}; jaraknya $uh/v_s
= 2\times 0.5/0.03 = \qty{33}{m}$. Bijinya: $2\times 20/0.8 =
\qty{50}{m}$ --- satuan yang lebih berat yang dilepaskan dari tempat
lebih tinggi dengan sayap menempuh jarak kira-kira sejauh itu juga.
\end{solution}

\begin{solution}{exo:b2:plant-life-cycles:7}
$200\times 40\times 64 = \num{512000}$ spora tiap ental; $5.1$
\omterm{def:b2:plant-life-cycles:fern}{protalus}; $0.1$ \omterm{def:b2:plant-life-cycles:cycles}{sporofit} muda tiap entalnya.
\end{solution}

\begin{solution}{exo:b2:plant-life-cycles:8}
Sebutir biji adalah \omterm{def:b2:plant-life-cycles:heterospory}{bakal biji} yang ditahan: \omterm{def:b2:plant-life-cycles:cycles}{gametofit} betinanya harus
tetap berada pada tetuanya, terkurung, dan disuapi, sedangkan \omterm{def:b2:plant-life-cycles:cycles}{gametofit}
jantannya harus mengembara menuju kepadanya. Hal itu menuntut dua macam
spora dengan dua nasib. Satu-satunya spora tumbuhan homospor harus
dilepaskan agar tumbuh menjadi \omterm{def:b2:plant-life-cycles:cycles}{gametofit} dwikelamin yang bebas;
menahannya akan menahan fungsi jantannya pula lalu memaksakan pembuahan
sendiri pada tiap tumbuhannya, dan tidak akan ada apa pun yang
mengembara.
\end{solution}

\begin{solution}{exo:b2:plant-life-cycles:9}
Sporanya: $\tfrac43\pi(1.5\times 10^{-5})^3\times 1100 =
\qty{1.6e-11}{kg}$, yaitu \qty{16}{ng}; energinya $0.5\times 1.6\times
10^{-8}\times 2\times 10^{4} = \qty{1.6e-4}{J}$. Bijinya: \qty{5}{mg},
$0.5\times 5\times 10^{-3}\times 2\times 10^{4} = \qty{50}{J}$ ---
$3\times 10^{5}$ kali lebih banyak. Sporanya hanya dapat menumbuhkan
satu sisik berfotosintesis yang kecil dan tidak apa-apa lagi sebelum
ada cahaya; bijinya dapat menumbuhkan akar dan pucuk di dalam gelap,
menunggu satu musim, serta selamat dari dimakan atau dikeringkan.
\end{solution}

\begin{solution}{exo:b2:plant-life-cycles:10}
$10^{11}/(10^{6}\times 20) = \num{5000}$ butir tiap meter kubik;
fluksnya $5000\times 10^{-7}\times 2 = \qty{1e-3}{}$ butir tiap sekon;
kira-kira \qty{1000}{s}, seperempat jam, tiap butirnya. Runjungnya
harus terbuka dan siap terima selama berjam-jam sampai berhari-hari,
dan tepat ketika runjung jantannya melepaskan \omterm{def:b2:plant-life-cycles:heterospory}{serbuk sari} ---
penyerbukannya terjadwal sampai ke pekannya.
\end{solution}

\begin{solution}{exo:b2:plant-life-cycles:11}
Kediploidan menutupi \omterm{def:b2:genome-diversification:mutation}{mutasi} \omterm{def:b2:meiosis-heredity:vocabulary}{resesif} yang merugikan; sebuah tubuh
\omterm{def:b2:meiosis-heredity:meiosis}{diploid} dengan jaringan pembuluh dapat besar dan berumur panjang, dan
ketinggiannya membantu penangkapan cahaya sekaligus \omterm{thm:b2:plant-life-cycles:dispersal}{penyebaran spora};
spora yang dibuat tinggi pada sebuah \omterm{def:b2:plant-life-cycles:cycles}{sporofit} tersebar di udara,
sedangkan gamet harus berenang. Keturunan \omterm{def:b2:meiosis-heredity:meiosis}{haploidnya} bertahan karena
\omterm{def:b2:meiosis-heredity:meiosis}{meiosis} dan pembuahan menuntut adanya fase \omterm{def:b2:meiosis-heredity:meiosis}{haploid}, dan butir \omterm{def:b2:plant-life-cycles:heterospory}{serbuk
sari} serta kantung lembaganya adalah tubuh minimal yang mengantarkan
gamet kepada satu sama lain lalu menyuapi lembaganya.
\end{solution}

\begin{solution}{exo:b2:plant-life-cycles:12}
Dengan membiakkan spora lalu mengikutinya di bawah mikroskop, ia dapat
menetapkan bahwa sebuah spora tumbuh menjadi tubuh kecil yang
mengusung organ kelamin, bahwa lembaganya muncul dari sel telur yang
telah dibuahi di dalam \omterm{def:b2:plant-life-cycles:moss}{arkegoniumnya}, bahwa tumbuhan pengusung sporanya
tumbuh darinya, dan bahwa organ yang sama ada dalam bentuk yang
menyusut di dalam \omterm{def:b2:plant-life-cycles:heterospory}{bakal biji} konifer: yakni pergiliran sebagai
rangkaian tubuh dan organ. Ia tidak dapat mengetahui apa yang
membedakan kedua keturunannya pada taraf sel. Cacah kromosom
menunjukkan bahwa kedua tubuhnya berbeda dengan faktor dua dalam cacah
kromosomnya, bahwa \omterm{def:b2:meiosis-heredity:meiosis}{meiosis} duduk pada pembentukan sporanya dan
pembuahannya pada sel telurnya, sehingga menjelaskan mengapa
pergilirannya wajib.
\end{solution}

\begin{solution}{pb:b2:plant-life-cycles:1}
\textbf{1.} $10\times 200\times 40\times 64 = 5.1\times 10^{6}$
spora.
\textbf{2.} $\tfrac43\pi(1.5\times 10^{-5})^3\times 1100 =
\qty{1.6e-11}{kg}$ tiap spora; panennya $8\times 10^{-5}$ kg, yaitu \qty{80}{mg}.
\textbf{3.} $v_s = 2\times 2.25\times 10^{-10}\times 1099\times
9.81/(1.62\times 10^{-4}) = \qty{0.030}{m/s}$.
\textbf{4.} $x = 1.5\times 0.5/0.030 = \qty{25}{m}$; cakram seluas
$\pi\times 25^2 = \qty{2000}{m^2}$; kira-kira \num{2600} spora tiap
meter persegi.
\textbf{5.} Arus udara naik sebesar beberapa sentimeter tiap sekon
melampaui $v_s$, sehingga spora yang terperangkap olakan tetap
melayang berjam-jam lalu menempuh ratusan kilometer; satu spora
mendirikan sebuah populasi kalau \omterm{def:b2:plant-life-cycles:fern}{protalusnya} sanggup membuahi diri
sendiri (ia mengusung kedua organnya), dan itulah sebabnya pulau
terpencil memiliki flora paku yang kaya.
\textbf{6.} Sporanya: $0.5\times 1.6\times 10^{-8}\times 2\times 10^{4} =
\qty{1.6e-4}{J}$; panennya: \qty{820}{J}; satu biji pinus: \qty{50}{J}
--- seluruh panen sporanya setara enam belas biji.
\textbf{7.} $5.1\times 10^{6}/2\times 10^{4} = 256$ \omterm{def:b2:plant-life-cycles:fern}{protalus}.
\textbf{8.} $1 - 0.8^{6} = 1 - 0.26 = 0.74$.
\textbf{9.} $256\times 0.74\times 0.5\times 0.02 = 1.9$ \omterm{def:b2:plant-life-cycles:fern}{tumbuhan paku}
dewasa tiap musim.
\textbf{10.} $30\times 1.9 \approx 57$: jauh lebih banyak daripada satu
pengganti sebuah populasi yang mantap, sehingga angka kebertahanannya
terlalu murah hati --- di habitat yang penuh, sebagian besar \omterm{def:b2:plant-life-cycles:cycles}{sporofit}
mudanya mati karena berdesakan dan ternaungi.
\textbf{11.} $10^{-3}/10^{-4} = \qty{10}{s}$. Banyak \omterm{def:b2:plant-life-cycles:fern}{tumbuhan paku}
mematangkan \omterm{def:b2:plant-life-cycles:moss}{anteridium} dan \omterm{def:b2:plant-life-cycles:moss}{arkegoniumnya} pada waktu yang berbeda, atau
melepaskan sebuah hormon (anteridiogen) yang menjadikan \omterm{def:b2:plant-life-cycles:fern}{protalus} muda
di sebelahnya jantan, sehingga spermanya datang dari individu lain.
\textbf{12.} Ia setebal satu sel, tak berakar, tak terlindungi,
membutuhkan air cair untuk pembuahan, dan hidup beberapa pekan: nyaris
seluruh kematian daurnya jatuh padanya (satu spora dari
tiap $2\times 10^{4}$ menjadi \omterm{def:b2:plant-life-cycles:fern}{protalus}; seperempat di antaranya tak
pernah melihat malam yang basah).
\textbf{13.} $v_s = 2\times 6.25\times 10^{-10}\times 599\times
9.81/(1.62\times 10^{-4}) = \qty{0.045}{m/s}$.
\textbf{14.} $3\times 15/0.045 = \qty{1000}{m}$.
\textbf{15.} $10^{11}/(2\times 10^{7}) = \num{5000}$ tiap meter kubik.
\textbf{16.} $5000\times 10^{-7}\times 2 = \qty{1e-3}{s^{-1}}$: 3,6 tiap
jam; yang pertama tiba sesudah kira-kira \qty{1000}{s}, yaitu
seperempat jam.
\textbf{17.} Lima kilometer searah angin, awannya telah menyebar ke
samping dan ke atas serta sebagian besar butirnya telah mengendap
(jangkauan \qty{1}{km} tiap \qty{15}{m} ketinggian), sehingga
kepekatannya beberapa orde besar lebih rendah dan sebuah \omterm{def:b2:plant-life-cycles:heterospory}{bakal biji}
mungkin menunggu berhari-hari untuk sebutir; \omterm{def:b2:plant-life-cycles:heterospory}{serbuk sari} pohonnya
sendiri kebanyakan memberi biji hasil penyerbukan sendiri, yang gugur.
Penyerbukan angin membutuhkan awan yang padat, maka rumpunnya.
\textbf{18.} $10^{14}$ butir bagi $2\times 10^{6}$ \omterm{def:b2:plant-life-cycles:heterospory}{bakal biji}: $5\times
10^{7}$ butir tiap \omterm{def:b2:plant-life-cycles:heterospory}{bakal biji}.
\textbf{19.} $2\times 10^{6}$ \omterm{def:b2:plant-life-cycles:heterospory}{bakal biji} $\times 0.6\times 0.8 = 9.6\times
10^{5}$ biji; \qty{4.8}{kg}; \qty{48}{MJ}.
\textbf{20.} $3\times 20/0.8 = \qty{75}{m}$: \omterm{def:b2:plant-life-cycles:heterospory}{serbuk sarinya} menempuh
satu kilometer, bijinya beberapa kali tinggi pohonnya; gennya
mengembara lewat \omterm{def:b2:plant-life-cycles:heterospory}{serbuk sari}, tumbuhannya lewat biji.
\textbf{21.} $9.6\times 10^{5}\times 0.05\times 0.01 = 480$ keturunan
dewasa, terhadap 57 milik \omterm{def:b2:plant-life-cycles:fern}{tumbuhan pakunya}.
\textbf{22.} Pinusnya: $4.8\times 10^{7}/480 = \qty{100}{kJ}$ tiap
keturunan dewasa; \omterm{def:b2:plant-life-cycles:fern}{tumbuhan pakunya}: $820/1.9 = \qty{430}{J}$ tiap
keturunan dewasa, dua ratus kali lebih sedikit.
\textbf{23.} Keduanya separuhnya bahan kering pada \qty{20}{kJ/g},
sehingga pada laju basal, katakanlah, \qty{0.5}{mW} tiap gram kering,
keduanya bertahan $10^{4}\,
\text{J/g}/(5\times 10^{-4}\,\text{W/g}) = 2\times 10^{7}$ s, yaitu
kira-kira delapan bulan: lamanya dormansi bukanlah yang dibeli energi
bijinya. Yang dibelinya adalah apa yang terjadi sesudah perkecambahan
--- sebuah akar dan sebuah pucuk yang dibangun di dalam gelap dari
cadangannya, berpekan-pekan sebelum kecambahnya menyuapi dirinya
sendiri --- dan, bersama kulitnya, perlindungan selama itu;
\qty{1.6e-4}{J} sporanya membangun beberapa sel yang harus segera
berfotosintesis.
\textbf{24.} Bijinya: pembuahan tanpa air; sebuah lembaga yang
terbekali, terlindungi, dan dorman yang mapan di tempat kering atau
bermusim; penyebaran lewat sayap, buah, dan hewan. Siasat \omterm{def:b2:plant-life-cycles:fern}{tumbuhan
pakunya} menang di habitat yang basah, ternaungi, dan mantap, serta
dalam menjajah tanah yang jauh atau yang baru terbuka, tempat propagul
yang murah, ringan, dan banyak lebih berarti daripada pembekalan.
\textbf{25.} \omterm{def:b2:plant-life-cycles:fern}{Tumbuhan pakunya}: $5.1\times 10^{6}$ spora tiap musim bagi
1,9 individu dewasa, yaitu $2.7\times 10^{6}$ spora tiap keturunan
dewasa, masing-masing \qty{430}{J}; pinusnya: $9.6\times 10^{5}$ biji
bagi 480 individu dewasa, yaitu \num{2000} biji tiap keturunan dewasa,
masing-masing \qty{100}{kJ}.
\end{solution}
